This paper engages with the empirical Bayes SAEM-MCMC procedure for decomposable Gaussian graphical models \parencite{donnet2012empirical}. We do not propose a new method. We assess an existing one and characterize the regimes in which its central decomposability assumption is, and is not, costly in practice.

Our main objective is to evaluate how the decomposability assumption distorts inference in practice. The assumption is convenient because it admits closed-form marginal likelihoods on the cliques and separators of the recovered graph; the cost of that convenience is what this paper measures. When the underlying true graph is not chordal, the estimated posterior concentrates on the set of minimal triangulations of $G^*$. The empirical implication is mild: meaningful structure is still recovered, and for moderate-$p$ problems decomposability is a reasonable assumption and a good compromise between tractability and performance.

We constructed the DIP as a way to exploit Stage-1's tractability further but found no accuracy improvements at the calibrated mixing weight (Section~\ref{sec:sim}; Appendix~\ref{sec:lambda}). A more promising direction is whether running Stage~1 first improves Stage~2's mixing time rather than its asymptotic recovery.

For performance, in low-dimensional settings, Bayesian methods perform as well, if not generally better, than frequentist methods, and allow for uncertainty quantification, which can be important in applied settings. Moreover, in settings where the partial correlation structure is suspected to be weak, Bayesian methods clearly outperformed frequentist approaches. Their biggest drawback is computational, as the sampling procedure over graphs quickly becomes very expensive with dimensionality. The two-stage procedure adds a further structural cost: Stage~1's MH proposals can be rejected (many iterations leave the graph unchanged) and its acceptance ratio reads marginal-likelihood ratios between HIW posteriors via clique- and separator-level Cholesky operations, while BDMCMC accepts every birth-death event by construction and uses the closed-form normalizing-constant approximation of \textcite{mohammadi2023accelerating}. The two-stage compute budget is therefore earned through the diagnostic information in its two-graph output, not through any improvement in recovery quality.

On the DDLY bank-volatility panel ($p = 106$), Stage~2 returns a non-decomposable graph with $93\%$ of its added edges breaking decomposability when injected into Stage~1, confirming on real data that the decomposability restriction is the binding constraint Stage~2 lifts.

Our main takeaway is that any GGM estimator must impose some structural assumption. Decomposable models are an attractive choice not only because their performance can be comparable to unrestricted methods, but because the structural consequences of the assumption are interpretable. The posterior concentrates on minimal triangulations; the practitioner can infer from the output whether the data reasonably accommodate a non-decomposable structure; and the tractability of the closed-form likelihood makes decomposable models a sensible default for moderate-$p$ problems.

The diagnostic potential of the two-stage output suggests a promising avenue for future research. We have suggestive evidence that comparing the Stage-1 graph $\hat G_D$ against the unrestricted Stage-2 graph $\hat G$ is informative about misspecification of the decomposability assumption. The diagnostic verdict itself appears robust to the choice of Stage-2 prior, even though the specific Stage-2 structure is not, which strengthens the case for treating the two-stage output as a diagnostic rather than as a unique non-decomposable point estimate when the verdict flags misspecification. A natural next step is to formalize this comparison as a Bayesian test of decomposability, using the posterior mass on the decomposable subspace, $\mathbb{P}(G \in \mathcal{D}_p \mid Y)$, as the test statistic, and the Bayes factor for $G \in \mathcal{D}_p$ versus $G \notin \mathcal{D}_p$ as the corresponding hypothesis test. In practice, $\mathbb{P}(G \in \mathcal{D}_p \mid Y)$ can be estimated as the fraction of post-burn-in Stage-2 samples that are decomposable.
